% ============================================================
% LAMPIRAN D
% ============================================================

\section*{LAMPIRAN D}
\addcontentsline{toc}{section}{LAMPIRAN D: Instrumen Evaluasi Ahli Paleografi}

\vspace{1em}

\begin{center}
{\Large\textbf{INSTRUMEN EVALUASI}}\\[0.5em]
{\Large\textbf{RUBRIK PENILAIAN AHLI PALEOGRAFI}}\\[0.5em]
{\large\textbf{Restorasi Dokumen Terdegradasi}}
\end{center}

\vspace{2em}

% ============================================================
% INFORMASI EVALUATOR
% ============================================================

\subsection*{Petunjuk Penggunaan Instrumen}

Rubrik ini dirancang untuk evaluasi kualitas restorasi dokumen historis terdegradasi oleh ahli paleografi atau restorasi arsip. Setiap kriteria dinilai dengan skala 1--5, kemudian diagregasi menggunakan metodologi yang dijelaskan di bawah. Instrumen ini telah divalidasi dalam studi percontohan dengan 15 dokumen ANRI (Mei 2024) oleh 2 ahli independen dengan reliabilitas tinggi (Cohen's $\kappa$ = 0,78).

\vspace{2em}

% ============================================================
% KRITERIA PENILAIAN
% ============================================================

\subsection*{Kriteria Penilaian}

Penilaian dilakukan berdasarkan 5 aspek utama dengan skala 1-5:
\begin{itemize}[leftmargin=*, itemsep=0pt]
    \item \textbf{Skor 5}: Sangat Baik (Kualitas optimal, sangat membantu keterbacaan)
    \item \textbf{Skor 4}: Baik (Kualitas baik, meningkatkan keterbacaan signifikan)
    \item \textbf{Skor 3}: Cukup (Kualitas memadai, ada peningkatan keterbacaan)
    \item \textbf{Skor 2}: Kurang (Kualitas kurang, peningkatan keterbacaan minimal)
    \item \textbf{Skor 1}: Sangat Kurang (Tidak ada peningkatan atau memperburuk)
\end{itemize}

\vspace{1em}

% ============================================================
% TABEL RUBRIK PENILAIAN
% ============================================================

\begin{longtable}{|p{0.5cm}|p{3.5cm}|p{7cm}|p{1.5cm}|}
\caption{Rubrik Penilaian Kualitas Restorasi Dokumen} \\
\hline
\textbf{No} & \textbf{Aspek Penilaian} & \textbf{Deskripsi} & \textbf{Skor} \\
\hline
\endfirsthead

\multicolumn{4}{c}{\textit{Lanjutan Tabel Rubrik Penilaian}} \\
\hline
\textbf{No} & \textbf{Aspek Penilaian} & \textbf{Deskripsi} & \textbf{Skor} \\
\hline
\endhead

\hline
\multicolumn{4}{r}{\textit{Bersambung ke halaman berikutnya}} \\
\endfoot

\hline
\endlastfoot

% ============================================================
% ISI RUBRIK
% ============================================================

\multicolumn{4}{|l|}{\textbf{A. KUALITAS VISUAL}} \\
\hline
1 & Kejernihan Piksel & Eliminasi noise tanpa merusak tinta & 1 2 3 4 5 \\
\hline
2 & Kesamaan Struktural & Preservasi bentuk dan proporsi huruf & 1 2 3 4 5 \\
\hline
3 & Kealamian Visual & Naturalitas hasil, tekstur kertas terjaga & 1 2 3 4 5 \\
\hline

\multicolumn{4}{|l|}{\textbf{B. AKURASI BINARISASI}} \\
\hline
4 & Pemisahan Teks-Latar & Deteksi teks lengkap, latar bersih & 1 2 3 4 5 \\
\hline
5 & Kontinuitas Goresan & Keutuhan garis dan sambungan antar huruf & 1 2 3 4 5 \\
\hline

\multicolumn{4}{|l|}{\textbf{C. KETERBACAAN KARAKTER}} \\
\hline
6 & Identifikasi Karakter & Huruf terbaca jelas tanpa ambiguitas & 1 2 3 4 5 \\
\hline
7 & Disambiguasi Huruf & Pembedaan huruf mirip (e/o, a/u, m/n, i/l) & 1 2 3 4 5 \\
\hline
8 & Detail Mikro & Ketajaman tanda diakritik dan baca & 1 2 3 4 5 \\
\hline

\multicolumn{4}{|l|}{\textbf{D. KETERBACAAN KATA}} \\
\hline
9 & Identifikasi Kata & Kata terbaca utuh, batas jelas & 1 2 3 4 5 \\
\hline
10 & Koherensi Spasi & Jarak antar karakter dan kata konsisten & 1 2 3 4 5 \\
\hline
11 & Aliran Teks & Kontinuitas teks dalam baris dan kalimat & 1 2 3 4 5 \\
\hline

\multicolumn{4}{|l|}{\textbf{E. PELESTARIAN HISTORIS}} \\
\hline
12 & Autentisitas Paleografi & Preservasi gaya tulisan periode & 1 2 3 4 5 \\
\hline
13 & Integritas Konten & Tidak ada perubahan informasi tekstual & 1 2 3 4 5 \\
\hline
14 & Kelayakan Transkripsi & Kemudahan transkripsi ahli paleografi & 1 2 3 4 5 \\
\hline
15 & Standar Arsip Digital & Kesesuaian standar preservasi UNESCO & 1 2 3 4 5 \\
\hline

\end{longtable}

\vspace{2em}

% ============================================================
% PERHITUNGAN SKOR
% ============================================================

\subsection*{Metodologi Agregasi Skor}

Rubrik 15 kriteria di atas diagregasi menjadi 4 kriteria untuk analisis makro:

\begin{table}[h]
\centering
\small
\begin{tabular}{|l|l|l|}
\hline
\textbf{Kriteria Agregat} & \textbf{Mapping dari Rubrik} & \textbf{Perhitungan} \\
\hline
Keterbacaan Teks & Kriteria 6--11 (C+D) & Rerata 6 kriteria \\
Autentisitas Paleografi & Kriteria 12--13 (E.1--2) & Rerata 2 kriteria \\
Minimasi Artefak & Kriteria 4--5 (B) & Rerata 2 kriteria (inverse) \\
Kegunaan Transkripsi & Kriteria 14--15 (E.3--4) & Rerata 2 kriteria \\
\hline
\multicolumn{3}{l}{\footnotesize Kriteria 1--3 (A) tidak dihitung dalam agregat (untuk validasi tambahan).} \\
\multicolumn{3}{l}{\footnotesize Konversi skala: Skor rubrik (1-5) $\times$ 0,8 = skor agregat (skala 1-4)} \\
\multicolumn{3}{l}{\footnotesize Skor keseluruhan: Rerata dari 4 kriteria agregat} \\
\end{tabular}
\end{table}

\subsection*{Validasi Instrumen: Hasil Studi Percontohan}

Rubrik di atas telah divalidasi dalam studi percontohan dengan 15 dokumen ANRI (Mei 2024) yang dievaluasi oleh 2 ahli independen. Hasil agregat menunjukkan reliabilitas instrumen:

\begin{table}[h]
\centering
\small
\begin{tabular}{|l|c|c|c|}
\hline
\textbf{Kriteria} & \textbf{Penilai 1} & \textbf{Penilai 2} & \textbf{Rerata} \\
\hline
Keterbacaan Teks & 3,47 & 3,33 & 3,40 \\
Autentisitas Paleografi & 3,60 & 3,53 & 3,57 \\
Minimasi Artefak (inv.) & 3,27 & 3,13 & 3,20 \\
Kegunaan Transkripsi & 3,53 & 3,47 & 3,50 \\
\hline
\textbf{Skor Keseluruhan} & \textbf{3,47} & \textbf{3,37} & \textbf{3,42} \\
\hline
\multicolumn{4}{l}{\footnotesize Skala: 1=Buruk, 2=Cukup, 3=Baik, 4=Sangat Baik} \\
\multicolumn{4}{l}{\footnotesize Cohen's kappa: $\kappa$ = 0,78 (kesepakatan substansial)} \\
\multicolumn{4}{l}{\footnotesize Spearman's $\rho$ = 0,91 (p$<$0,001) untuk ranking dokumen} \\
\multicolumn{4}{l}{\footnotesize 95\% CI untuk rerata: [3,28, 3,56]} \\
\end{tabular}
\end{table}

\textit{Sumber}: Lihat Chapter 5 untuk analisis lengkap hasil evaluasi. Distribusi dokumen: kontras tinggi (n=5, $\bar{x}$=3,75), sedang (n=8, $\bar{x}$=3,38), rendah (n=2, $\bar{x}$=2,63). Konsistensi antarpenilai ($\kappa$=0,78) mengonfirmasi validitas instrumen untuk evaluasi dokumen historis terdegradasi.

\vspace{2em}

% ============================================================
% CATATAN EVALUATOR
% ============================================================

\subsection*{Ringkasan Umpan Balik Ahli}

\textbf{Kelebihan Utama:}
\begin{itemize}[leftmargin=*, itemsep=2pt]
    \item Preservasi karakteristik skrip Belanda abad ke-17 (ligatur, \textit{aspect ratio})
    \item Pembersihan \textit{foxing} efektif tanpa korbankan detail autentisitas
    \item Detail mikroskopis goresan terjaga untuk analisis paleografi
    \item Sesuai untuk digitisasi massal dengan koreksi manual minimal
\end{itemize}

\textbf{Area Perbaikan:}
\begin{itemize}[leftmargin=*, itemsep=2pt]
    \item \textit{Slight blur} pada region tinta \textit{iron gall} terkorosi (2--3 kasus sedang)
    \item Derau residual minimal pada \textit{bleed-through} berat
    \item Disarankan verifikasi ahli untuk dokumen kontras <30\%
\end{itemize}

\textbf{Rekomendasi:}
Pendekatan konservatif model yang menghindari \textit{creative reconstruction} sangat diapresiasi untuk preservasi integritas historis. Cocok untuk praproses \textit{batch} pipeline HTR dengan supervisi minimal.

\vspace{2em}

\subsection*{Validitas dan Reliabilitas}

Evaluasi dilakukan oleh 2 evaluator independen dengan kesepakatan substansial (Cohen's $\kappa$ = 0,78). Korelasi Spearman antarpenilai: $\rho$ = 0,91 (p$<$0,001). Hasil mendukung validitas \textit{pilot study} untuk eksplorasi aplikasi arsip praktis meskipun ukuran sampel terbatas (n=15).

